% =====================================================================
%  Chapter 3 (Proposed Design) -- PART B: Design Solution
%  This file CONTINUES the chapter opened by Part A. It does NOT call
%  \chapter. The first section carries the chapter-continuation label.
% =====================================================================

\section{Design Alternatives and Justification}
\label{ch:design_solution}

Part~A of this chapter set out \emph{what} SmartLoad has to do: the functional
and non-functional requirements, the four-context concerns, and the risk
register that any design must answer to. This second part explains \emph{how}
we chose the design. For each important decision, we listed several candidate
designs, weighed them carefully against the requirements and risks, and then
picked one. The same tension runs through every decision: the system has to be
smart enough to adapt to changing load, but that smartness must never sit on
the critical path in a way that can take traffic down. The rule that settles
this tension is \emph{graceful degradation by default} --- if any adaptive
component is lost, SmartLoad drops back to classical routing, never to an
outage. We judge each alternative below first on whether it respects that rule,
and only then on its performance, cost, and how easy it is to extend.

\subsection{Routing Intelligence}
\label{dsb:alt-routing}

Routing is the decision the data plane makes on every request: given a pool of
backends, which one gets the next connection. We looked at four candidate
designs.

\textbf{(a) Static round-robin.} The classical default: cycle through the pool
in order and give each backend an equal share. It is deterministic, has no
dependencies, and is trivially correct, which is exactly why we keep it as the
deterministic floor. The limitation is that it is blind --- it cannot react to
a slow or unhealthy backend, and it wastes capacity when the backends are not
all the same size.

\textbf{(b) Heuristic least-connections.} Send each request to the backend with
the fewest in-flight connections. This is a small step up: it reacts to load
imbalance with no model and no training. But it reacts to only one coarse
signal (connection count), it knows nothing about latency or error rate, and it
cannot look ahead --- it is purely reactive.

\textbf{(c) A learned routing policy.} Train a reinforcement-learning policy
\cite{sutton2018rl} that watches per-backend telemetry (latency, queue depth,
health) and produces a ranking that the data plane turns into upstream weights.
In principle this captures patterns that neither heuristic can --- correlated
slowdowns, backends with different capacity, anticipatory shifting. The
trade-offs are just as clear: it needs offline training
\cite{schulman2017ppo,raffin2021sb3}, its behaviour is harder to explain, and a
badly trained or stale policy could route traffic poorly. A learned policy that
is allowed to authorise routing on its own does not fit the degradation rule.

\textbf{(d) A hybrid learned policy with a deterministic fallback (chosen).}
Run the learned policy, but only ever as an \emph{advisory} input bounded by two
deterministic guards: an anomaly check that strips unhealthy backends out of the
action space before the policy even sees them, and a classical floor that takes
over whenever the policy is disabled, warming up, stale, or overridden. The
learned policy proposes; deterministic logic disposes. This keeps the upside of
(c) --- adaptivity --- while structurally bounding the downside, because the
worst case degrades to (a). This is the design SmartLoad adopts, and
\cref{dsb:sec-decision} shows in numbers why it is the clear winner.

\subsection{Where Routing Decisions Live}
\label{dsb:alt-locus}

A second, separate decision is \emph{where} the routing intelligence physically
runs relative to the data plane.

\textbf{(a) An in-process load-balancer module.} Put the decision logic inside
the data plane itself --- for example as an NGINX \cite{nginx} module or Lua
handler that computes weights per request. This is the lowest-latency option,
but it breaks the separation of decision and data planes: model code, database
queries, and training dependencies would now live in the request path, where a
fault stalls traffic. It also ties the intelligence to one specific load
balancer.

\textbf{(b) An external control plane that writes upstream configuration
(chosen).} Keep the data plane (NGINX) free of intelligence, and place every
decision service off the request path. A single bridging process --- the
\emph{LB sidecar} --- subscribes to the control bus and regenerates the NGINX
upstream configuration, then applies it with a graceful reload. Decisions are
computed asynchronously and consumed as pre-computed signals; the request path
reads only its in-memory effective configuration. A fault in any decision
service stops \emph{adaptation}, not \emph{traffic}. This is SmartLoad's
architecture.

\textbf{(c) A full service mesh.} Adopt a mesh such as Istio \cite{istio} and
express routing as mesh policy enforced by per-pod sidecars. A mesh is powerful
and proven in production, but it is too heavy for this problem: it brings a
control plane, a data plane of injected proxies, and an operational learning
curve far larger than a thin middleware needs, and it does not give us the
telemetry-driven forecasting and learned-routing loop that is SmartLoad's main
contribution. The mesh would become the platform rather than a component inside
it, which goes against the middleware-first principle.

\subsection{Forecasting}
\label{dsb:alt-forecast}

Forecasting predicts the short-horizon request rate so the autoscaler can
provision ahead of demand rather than after it.

\textbf{(a) A moving-average baseline.} Predict the next horizon as the mean of
a recent window, with one standard deviation as a spread indicator. It is exact,
cheap, and always produces an answer, which is why it ships as the forecasting
baseline. What it cannot do is capture trend or seasonality.

\textbf{(b) ARIMA.} Classical Box--Jenkins time-series modelling
\cite{box2015timeseries,hyndman2021forecasting} fits autoregressive and
moving-average terms and gives a principled prediction interval. It captures
trend and short-horizon momentum that a flat mean cannot, in exchange for an
offline-trained artefact and heavier inference. SmartLoad ships an ARIMA engine
as the trained replacement for the baseline.

\textbf{(c) Heavier deep models.} Recurrent or attention-based sequence models
can in principle model long, complex seasonalities. For a five-minute scaling
horizon on operational request-rate data, we judged the extra accuracy not worth
the training cost, the artefact size, and the inference latency --- the
autoscaler needs a useful number quickly, not a slightly better number slowly.
We record deep models as future work, not a current choice.

\subsection{Anomaly Detection}
\label{dsb:alt-anomaly}

Anomaly detection flags backends that are degraded or unhealthy so they can be
excluded from routing before the learned policy ever sees them.

\textbf{(a) A static threshold rule.} Flag a backend when its current latency
goes above a multiple of its own rolling mean (or its error rate crosses a
bound). It is interpretable, needs no model, and always produces a
classification --- the properties that make it the anomaly baseline and the
canonical fallback.

\textbf{(b) Isolation Forest.} An unsupervised tree ensemble
\cite{liu2008isolationforest} that isolates outliers in feature space without a
labelled training set. It can catch multivariate anomalies the threshold rule
misses, and it is light enough to score every cycle. SmartLoad ships an
Isolation Forest engine as the trained replacement, with the threshold rule as
its guaranteed fallback.

\textbf{(c) A deep autoencoder.} Reconstruction-error detectors such as those
benchmarked on the Server Machine Dataset \cite{su2019omnianomaly} can model
rich temporal structure, but they need a lot more data and tuning, and their
opacity works against the transparency the operator surface needs for a
safety-critical exclusion. We record them as future work.

\subsection{Messaging Between Services}
\label{dsb:alt-messaging}

The decision services have to hand their signals to the actuating services.

\textbf{(a) Synchronous HTTP.} Each service calls the next directly over HTTP.
This is simple and easy to trace, but it couples producers to consumers: a
publisher blocks on a slow or absent subscriber, and adding a new consumer (for
example the operator UI's live feed) means editing the producer. For a fan-out
of advisory signals, synchronous request/response is the wrong shape.

\textbf{(b) A Redis pub/sub control bus (chosen).} Decision services publish
typed envelopes onto named channels, and any number of consumers subscribe
\cite{carlson2013redis}. Publishers never block on subscribers, new consumers
attach without touching producers, and the fire-and-forget behaviour fits the
nature of the signals --- a missed advisory message is harmless because the next
cycle re-publishes, and the only durable surface (the audit hypertable) is
written separately. This decoupling is what lets a decision service fail without
stalling any other. SmartLoad uses five such channels, detailed in
\cref{dsb:sec-bus}.

\section{Decision Analysis}
\label{dsb:sec-decision}

To make the routing-architecture decision something we could defend rather than
just assert, we applied a weighted decision matrix in the Pugh tradition
\cite{pugh1991totaldesign}. We drew six criteria straight from the requirements
and risks of Part~A and gave each a weight that reflects how much it matters for
a safety-critical middleware. We scored each candidate routing design on a
one-to-five scale per criterion (five is best); the weighted total is the sum of
weight~$\times$~score across the criteria.

The criteria, with the reasoning behind each weight, are: \emph{Safety} (weight
$0.30$) --- the dominant concern, capturing whether a failure or mistake in the
routing layer can take traffic down; \emph{Latency overhead} (weight $0.15$)
--- the cost the design adds to the request path; \emph{Operational weight}
(weight $0.15$) --- the effort to deploy, run, and reason about the design;
\emph{Extensibility} (weight $0.15$) --- how easily a better algorithm can be
dropped in later; \emph{Transparency} (weight $0.15$) --- whether an operator
can reconstruct why a routing decision was made; and \emph{Cost} (weight
$0.10$) --- compute, training, and dependency footprint. The weights sum to one.
\Cref{dsb:tab:decision} records the scoring.

\begin{table}[t]
\centering
\caption{Weighted decision matrix for the routing-architecture decision. Scores
are on a 1--5 scale (5 is best); the weighted total is $\sum_i w_i s_i$. The
hybrid learned-with-deterministic-fallback design maximises the weighted total
once the dominant safety weight is taken into account.}
\label{dsb:tab:decision}
\small
\begin{tabular}{l c cccc}
\toprule
 & & \multicolumn{4}{c}{Candidate routing design (score 1--5)} \\
\cmidrule(lr){3-6}
Criterion & Weight & Round- & Least- & Learned & Hybrid \\
          &        & robin  & conn.  & only    & (chosen) \\
\midrule
Safety              & 0.30 & 5 & 4 & 2 & 5 \\
Latency overhead    & 0.15 & 5 & 4 & 3 & 4 \\
Operational weight  & 0.15 & 5 & 4 & 2 & 3 \\
Extensibility       & 0.15 & 1 & 2 & 5 & 5 \\
Transparency        & 0.15 & 5 & 4 & 2 & 4 \\
Cost                & 0.10 & 5 & 4 & 2 & 3 \\
\midrule
\textbf{Weighted total} & \textbf{1.00} & \textbf{4.40} & \textbf{3.70} & \textbf{2.70} & \textbf{4.20} \\
\bottomrule
\end{tabular}
\end{table}

The matrix is clear about a tension it exposes. Pure round-robin scores highest
in raw weighted terms ($4.40$) precisely because it is safe, cheap, and
transparent --- but it scores the floor on extensibility, and on its own it
cannot meet the adaptive-routing requirement from Part~A at all: it has no way to
react to a slow backend or to a learned signal. So it is not a candidate for the
\emph{whole} design; it is the candidate for the \emph{fallback}. The key
observation is that the hybrid design ($4.20$) scores within a whisker of
round-robin on every safety-related criterion \emph{while} earning the top
extensibility score, because its worst case \emph{is} round-robin. In other
words, the hybrid does not trade safety for adaptivity --- it layers adaptivity
on top of a round-robin floor it can always fall back to. The learned-only
design ($2.70$) loses ground on every axis that matters: it gives up safety,
transparency, and operational simplicity for an adaptivity the hybrid also
delivers, but without the structural guard.

Reading the matrix against Part~A closes the argument. The dominant risk in the
register was that an adaptive routing layer could degrade availability; the
$0.30$ safety weight encodes that risk directly, and only the round-robin and
hybrid designs score $5$ there. Of those two, only the hybrid also satisfies the
adaptive-routing and replaceable-engine requirements. The selected design is
therefore the external control plane (\cref{dsb:alt-locus}) carrying a hybrid
learned-with-deterministic-fallback routing policy (\cref{dsb:alt-routing}): the
locus choice keeps intelligence off the request path, and the routing choice
keeps a deterministic floor underneath it. The rest of this chapter describes
that selected design in detail.

\section{Description of the Selected Design}
\label{dsb:sec-design}

SmartLoad is organised as a set of cooperating containerised services split
across three planes. The \emph{data plane} carries client traffic; the
\emph{decision plane} computes adaptive signals off the request path; and the
\emph{control plane} turns those signals into concrete actuations on the data
plane. The subsections below build the design up one view at a time --- from the
outside-in system context, through the plane split and the closed control loop,
down to the control bus, the engine abstraction, and the layered safety model
that makes the whole thing trustworthy.

\subsection{System Context}
\label{dsb:sec-context}

At the coarsest level, SmartLoad is middleware that sits between client traffic
and a pool of backend services. \Cref{fig:dsb:context} shows the context.
Clients send HTTP requests to the NGINX data plane \cite{nginx}, which proxies
them to a backend in the pool and returns the response; from the client's point
of view SmartLoad is an ordinary reverse proxy. Beside the data plane sits the
control plane, which never touches a request directly: it observes the pool
through telemetry and steers the data plane by rewriting its upstream
configuration. The operator works with the control plane through two read/write
surfaces --- a policy interface (to set bounds and the safe-mode kill switch)
and an operator UI plus Grafana dashboards \cite{grafana} (to see what every
engine is doing). The backend pool is the customer's application; SmartLoad
treats it as an opaque set of interchangeable endpoints.

\begin{figure}[t]
\centering
\begin{tikzpicture}[node distance=8mm and 16mm]
  \node[slnode] (client) {\textbf{Clients}\\HTTP/1.1};
  \node[slnode,right=of client] (nginx) {\textbf{SmartLoad}\\NGINX\\data plane};
  \node[slnode,right=of nginx] (pool) {\textbf{Backend}\\service pool\\($\times N$)};
  \node[slnode,below=14mm of nginx] (cp) {\textbf{Control plane}\\(decision +\\actuation,\\off-path)};
  \node[slnode,below=10mm of client] (op) {\textbf{Operator}\\policy +\\UI / Grafana};
  % request path
  \draw[slflow] (client) -- node[above,font=\footnotesize]{requests} (nginx);
  \draw[slflow] (nginx) -- node[above,font=\footnotesize]{proxy} (pool);
  \draw[slflow,dashed] (pool.south west) to[bend left=10] node[below,font=\footnotesize]{responses} (client.south east);
  % control linkage
  \draw[slbus] (cp) -- node[right,font=\footnotesize]{rewrite upstream} (nginx);
  \draw[slbus,dashed] (pool.south) |- node[pos=0.75,below,font=\footnotesize]{telemetry} (cp.east);
  \draw[slflow] (op) -- node[left,font=\footnotesize]{steer / inspect} (cp);
\end{tikzpicture}
\caption{System context. SmartLoad is a reverse proxy on the request path
(clients $\rightarrow$ NGINX $\rightarrow$ backend pool $\rightarrow$ clients).
The control plane sits beside the data plane: it observes the pool through
telemetry and steers routing by rewriting the NGINX upstream configuration,
never by touching a live request. The operator steers and inspects the control
plane.}
\label{fig:dsb:context}
\end{figure}

\subsection{Three-Plane Architecture}
\label{dsb:sec-planes}

\Cref{fig:dsb:planes} opens the box. The \emph{data plane} (NGINX and the
backend pool) is the only surface that handles live requests. The
\emph{decision plane} holds the four adaptive services --- anomaly detector,
forecaster, learned routing engine, and autoscaler --- which read aggregated
telemetry from TimescaleDB \cite{timescaledb} and emit signals; none of them
ever proxies traffic. The \emph{control plane} holds the policy manager, the
Redis control bus \cite{carlson2013redis}, and the LB sidecar that bridges back
to the data plane. Telemetry flows \emph{up} from the data plane (NGINX access
logs are tailed, shipped over OpenTelemetry \cite{opentelemetry}, and persisted
to TimescaleDB), and decisions flow \emph{down} (signals cross the bus and the
sidecar rewrites the upstream configuration). The key structural property,
visible in the figure, is that exactly one process --- the LB sidecar ---
crosses the plane boundary downward, and it does so by writing configuration,
not by taking part in request handling. Every adaptive service can fail without
stalling a single request.

\begin{figure}[t]
\centering
\begin{tikzpicture}[node distance=6mm and 10mm]
  % --- data plane ---
  \node[slnode] (nginx) {NGINX};
  \node[slnode,right=of nginx] (pool) {Backend pool};
  \node[slplane,fit=(nginx)(pool),label={[font=\footnotesize\bfseries]above:DATA PLANE}] (dp) {};
  % --- decision plane ---
  \node[slnode,below=16mm of nginx] (anom) {Anomaly\\detector};
  \node[slnode,right=of anom] (fore) {Forecaster};
  \node[slnode,right=of fore] (rl) {Routing\\engine};
  \node[slnode,right=of rl] (as) {Autoscaler};
  \node[slplane,fit=(anom)(fore)(rl)(as),label={[font=\footnotesize\bfseries]above:DECISION PLANE}] (decp) {};
  % --- control plane ---
  \node[slnode,below=16mm of anom] (pm) {Policy\\manager};
  \node[slnode,right=of pm] (bus) {Redis\\control bus};
  \node[slnode,right=of bus] (db) {TimescaleDB};
  \node[slnode,right=of db] (sc) {LB sidecar};
  \node[slplane,fit=(pm)(bus)(db)(sc),label={[font=\footnotesize\bfseries]above:CONTROL PLANE}] (cp) {};
  % --- flows ---
  \draw[slflow] (nginx) -- (pool);
  \draw[slbus,dashed] (dp.south) -- node[right,font=\footnotesize]{telemetry up} (decp.north);
  \draw[slflow] (decp.south) -- node[right,font=\footnotesize]{signals} (bus.north);
  \draw[slflow] (db.north) -- node[left,font=\footnotesize]{metrics} (decp.south);
  \draw[slbus] (sc.north) to[bend right=18] node[right,font=\footnotesize]{decisions down} (pool.south);
\end{tikzpicture}
\caption{Three-plane architecture. The data plane (top) serves traffic; the
decision plane (middle) reads telemetry from TimescaleDB and emits adaptive
signals; the control plane (bottom) owns policy, the Redis bus, durable
storage, and the LB sidecar. Telemetry flows up and decisions flow down. The LB
sidecar is the only process that crosses the boundary into the data plane, and
it does so by rewriting configuration, never by handling a request.}
\label{fig:dsb:planes}
\end{figure}

\subsection{The Adaptive Control Loop}
\label{dsb:sec-mapek}

The planes cooperate as a continuous closed-loop controller in the MAPE-K style
of autonomic computing \cite{kephart2003autonomic,ibm2005mapek}:
\emph{Monitor}, \emph{Analyze}, \emph{Plan}, \emph{Execute}, over shared
\emph{Knowledge}. \Cref{fig:dsb:mapek} maps SmartLoad onto this loop.
\emph{Monitor} is the telemetry pipeline that lands per-request metrics in
TimescaleDB. \emph{Analyze} is the anomaly detector, scoring each backend's
health. \emph{Plan} is the combination of the forecaster, the learned routing
engine, and the autoscaler, each turning analysis into a proposed action.
\emph{Execute} is the LB sidecar, which merges the proposals under policy and
applies them to NGINX, together with the autoscaler's pool changes. The shared
\emph{Knowledge} is the operating policy and the durable telemetry and audit
store. After each actuation, new telemetry reflects the post-action state and
the loop closes.

A representational note keeps the figure precise about what flows where.
\emph{Forecasts drive scaling only; they never enter the routing decision.}
A demand prediction adjusts how many backends exist, but a prediction error
must never be able to move a single per-request routing choice. This separation
of forecast (capacity) from routing (per-request selection) is a deliberate
safety boundary, not an accident of wiring.

\begin{figure}[t]
\centering
\begin{tikzpicture}[node distance=10mm and 18mm]
  \node[slnode] (mon) {\textbf{Monitor}\\telemetry pipeline\\$\rightarrow$ TimescaleDB};
  \node[slnode,right=of mon] (ana) {\textbf{Analyze}\\anomaly detector\\(health scoring)};
  \node[slnode,below=14mm of ana] (plan) {\textbf{Plan}\\forecast + routing\\engine + autoscaler};
  \node[slnode,left=of plan] (exec) {\textbf{Execute}\\LB sidecar\\(merge + apply)};
  \node[slnode,below=14mm of mon,xshift=14mm] (know) {\textbf{Knowledge}\\policy + TimescaleDB\\audit store};
  \draw[slflow] (mon) -- (ana);
  \draw[slflow] (ana) -- (plan);
  \draw[slflow] (plan) -- (exec);
  \draw[slflow] (exec.south west) to[bend right=20] node[left,font=\footnotesize]{new state} (mon.south west);
  \draw[slbus,dashed] (know.north) -- (ana.south);
  \draw[slbus,dashed] (know) -- (plan);
\end{tikzpicture}
\caption{The adaptive control loop, in the MAPE-K idiom. Monitor (telemetry)
$\rightarrow$ Analyze (anomaly) $\rightarrow$ Plan (forecast, routing engine,
autoscaler) $\rightarrow$ Execute (LB sidecar), over shared Knowledge (policy
and the durable audit store). Each actuation produces new telemetry, closing the
loop. Forecasts influence the Plan stage's scaling output only --- never the
routing decision.}
\label{fig:dsb:mapek}
\end{figure}

\subsection{Control Bus}
\label{dsb:sec-bus}

The decision and control planes communicate over a Redis pub/sub bus
\cite{carlson2013redis} carrying five named channels of typed envelopes. The
choice of pub/sub over synchronous HTTP (\cref{dsb:alt-messaging}) is what lets
a publisher fail without stalling a consumer, and a new consumer attach without
editing a producer. \Cref{fig:dsb:bus} shows the publishers and subscribers of
each channel. \texttt{smartload.policy} is published by the policy manager and
read by every service that needs the operating policy; the remaining four carry
decision signals: \texttt{smartload.anomaly} (anomaly detector $\rightarrow$ LB
sidecar and operator UI), \texttt{smartload.forecast} (forecaster $\rightarrow$
autoscaler), \texttt{smartload.routing} (routing engine $\rightarrow$ LB
sidecar), and \texttt{smartload.scale} (autoscaler $\rightarrow$ LB sidecar and
operator UI). The bus is fire-and-forget by design: a missed message is harmless
because the next cycle re-publishes, and the only durable record is the audit
hypertable, written separately.

\begin{figure}[t]
\centering
\begin{tikzpicture}[node distance=5mm and 22mm]
  % publishers (left)
  \node[slnode] (pm) {Policy\\manager};
  \node[slnode,below=of pm] (ad) {Anomaly\\detector};
  \node[slnode,below=of ad] (fe) {Forecaster};
  \node[slnode,below=of fe] (rl) {Routing\\engine};
  \node[slnode,below=of rl] (as) {Autoscaler};
  % channels (middle)
  \node[slnode,right=of pm,fill=slbox] (c1) {\texttt{smartload.policy}};
  \node[slnode,below=of c1] (c2) {\texttt{smartload.anomaly}};
  \node[slnode,below=of c2] (c3) {\texttt{smartload.forecast}};
  \node[slnode,below=of c3] (c4) {\texttt{smartload.routing}};
  \node[slnode,below=of c4] (c5) {\texttt{smartload.scale}};
  % subscribers (right)
  \node[slnode,right=of c1] (sall) {all services};
  \node[slnode,right=of c3] (sc) {LB sidecar};
  \node[slnode,below=of sc] (oui) {operator UI};
  % publish edges
  \draw[slbus] (pm) -- (c1);
  \draw[slbus] (ad) -- (c2);
  \draw[slbus] (fe) -- (c3);
  \draw[slbus] (rl) -- (c4);
  \draw[slbus] (as) -- (c5);
  % subscribe edges
  \draw[slflow] (c1) -- (sall);
  \draw[slflow] (c2) -- (sc);
  \draw[slflow] (c3) -- (as.east);
  \draw[slflow] (c4) -- (sc);
  \draw[slflow] (c5) -- (sc);
  \draw[slflow] (c2) -- (oui);
  \draw[slflow] (c5) -- (oui);
\end{tikzpicture}
\caption{The five-channel control bus. Publishers (left) emit typed envelopes
onto named Redis channels (centre); subscribers (right) consume them.
\texttt{smartload.policy} fans out to every service; the four decision channels
route to the LB sidecar (and, for anomaly and scale, the operator UI's live
feed). The forecaster's output goes only to the autoscaler --- forecasts never
reach the LB sidecar, keeping prediction out of the routing path.}
\label{fig:dsb:bus}
\end{figure}

\subsection{Pluggable Engine Contract}
\label{dsb:sec-engine}

The replaceable-engine principle is realised as a strict plugin contract shared
by the three decision services that host swappable algorithms (anomaly
detection, forecasting, and learned routing). \Cref{fig:dsb:engine} shows the
abstraction. Three rules govern every plugin. First, \emph{plugins are pure}: an
engine is a function from a structured input dataclass to a structured output
dataclass --- it opens no sockets, queries no database, touches no bus, and reads
no environment. Every side effect lives in the surrounding \emph{run loop},
which owns the telemetry query, the envelope construction, and the publish.
Second, \emph{the run loop owns the envelope}: a plugin never builds a wire
message; the loop converts the plugin's output dataclass into the published
payload. Third, \emph{the baseline never fails}: each service ships one baseline
engine with no model artefact and a constructor that cannot raise, and the run
loop installs it automatically whenever the requested engine fails to load.

The result is a clean three-state engine: a deterministic \emph{baseline}, a
\emph{trained} replacement that drops in as one folder behind the same
interface, and an \emph{automatic fallback} from the trained engine to the
baseline on any load failure. Swapping the threshold rule for an Isolation
Forest \cite{liu2008isolationforest}, or the moving average for ARIMA
\cite{box2015timeseries}, is a change inside one service and one model file ---
the LB sidecar and autoscaler, which depend on the message contract rather than
the engine, never notice.

\begin{figure}[t]
\centering
\begin{tikzpicture}[node distance=8mm and 14mm]
  \node[slnode] (db) {TimescaleDB\\rows};
  \node[slnode,right=of db] (loop) {\textbf{Run loop}\\(owns I/O +\\the envelope)};
  \node[slnode,above right=8mm and 16mm of loop] (trained) {\textbf{Trained}\\engine\\(model artefact)};
  \node[slnode,below right=8mm and 16mm of loop] (base) {\textbf{Baseline}\\engine\\(no artefact,\\cannot fail)};
  \node[slnode,right=34mm of loop] (bus) {Redis\\channel};
  \draw[slflow] (db) -- node[above,font=\footnotesize]{query} (loop);
  \draw[slflow] (loop) -- node[above left,font=\footnotesize]{input dataclass} (trained);
  \draw[slflow,dashed] (trained) -- node[right,font=\footnotesize]{load fails} (base);
  \draw[slflow] (base) -- node[below right,font=\footnotesize]{output dataclass} (loop);
  \draw[slflow] (loop) -- node[above,font=\footnotesize]{envelope} (bus);
\end{tikzpicture}
\caption{The pluggable engine contract. The run loop owns all input/output ---
it queries TimescaleDB, hands the engine a pure input dataclass, receives a pure
output dataclass, and builds the published envelope. The engine is one of two
interchangeable implementations behind the same interface: a trained engine
(with a model artefact) or the baseline (no artefact, cannot fail). If the
trained engine fails to load, the run loop installs the baseline automatically.}
\label{fig:dsb:engine}
\end{figure}

\paragraph{Representational models.} Three small formulae pin down the behaviour
the figures describe. The learned routing engine emits a per-backend score $s_b$
over the surviving (non-excluded) backends $\mathcal{B}$; the LB sidecar turns
scores into integer NGINX upstream weights by normalising over the pool,
\begin{equation}
  w_b \;=\; \max\!\left(1,\ \left\lceil W_{\max}\cdot
  \frac{s_b}{\sum_{b'\in\mathcal{B}} s_{b'}}\right\rceil\right),
  \qquad b\in\mathcal{B},
  \label{dsb:eq:weights}
\end{equation}
so that every surviving backend keeps a floor weight of one (no backend is ever
starved) and relative preference follows the learned ranking. The autoscaler
converts a forecast into a target pool size by dividing predicted demand by
per-instance capacity,
\begin{equation}
  n^{\star} \;=\; \operatorname{clip}\!\left(
  \left\lceil \frac{\widehat{\mathrm{rps}}}{c}\right\rceil,\ n_{\min},\ n_{\max}\right),
  \label{dsb:eq:scale}
\end{equation}
where $\widehat{\mathrm{rps}}$ is the predicted request rate, $c$ is the
per-instance capacity, and $[n_{\min},n_{\max}]$ are the policy bounds. Finally,
the threshold anomaly rule flags a backend whenever its current latency goes
above a policy multiple of its own rolling mean,
\begin{equation}
  \texttt{unhealthy} \;\Longleftrightarrow\;
  \ell_{\text{cur}} \;>\; m\cdot \bar{\ell}_{\text{roll}},
  \label{dsb:eq:anomaly}
\end{equation}
with multiplier $m$ supplied by the operating policy. Each formula is computed
off the request path; the data plane consumes only the resulting weights.

\subsection{Three-Tier Safety Model}
\label{dsb:sec-safety}

The reason a learned policy can be trusted in the request path at all is that
three independent override tiers sit between it and live traffic, and \emph{all
three must agree} before any learned weight is applied. \Cref{fig:dsb:safety}
shows the layered gate. The first tier is the operating policy's
\texttt{safe\_mode} flag: when an operator sets it, the published routing mode is
forced to \emph{shadow}, weights reset to equal, and the data plane falls back to
round-robin --- a single human-controlled kill switch. The second tier is the
\texttt{RL\_MODE} deployment pin: even with \texttt{safe\_mode} off, the mode
stays \emph{shadow} unless the deployment explicitly pins
\texttt{RL\_MODE=active}, so a freshly deployed learned policy publishes its
inferences without actuating until an operator deliberately promotes it. The
third tier is the LB sidecar's own active gate: only when \texttt{safe\_mode} is
off, \texttt{RL\_MODE} is \texttt{active}, \emph{and} the policy itself returns
an \emph{active} action does the sidecar apply learned weights. The default
across all three tiers is the safe state; the trained policy alone can never
escalate routing past shadow.

This layered gate sits inside a strict routing hierarchy that the LB sidecar
enforces in order: anomaly exclusion first (an unhealthy backend is removed from
the action space and the learned policy never sees it), then policy bounds and
\texttt{safe\_mode}, then learned weights, then the classical floor. Anomaly
constraints outrank everything --- \texttt{safe\_mode} can silence the learned
policy but can never re-admit a backend that anomaly excluded --- because routing
to a known-bad backend has unbounded cost while the cost of conservatively
excluding a borderline one is bounded.

\begin{figure}[t]
\centering
\begin{tikzpicture}[node distance=9mm and 16mm]
  \node[slnode] (policy) {\textbf{Tier 1}\\policy \texttt{safe\_mode}\\(operator kill switch)};
  \node[slnode,right=of policy] (pin) {\textbf{Tier 2}\\\texttt{RL\_MODE} pin\\(deployment gate)};
  \node[slnode,right=of pin] (gate) {\textbf{Tier 3}\\LB sidecar active gate\\(policy returns active)};
  \node[slnode,below=12mm of pin,fill=slbox] (apply) {apply learned weights};
  \node[slnode,left=of apply] (shadow) {force \emph{shadow}\\$\rightarrow$ round-robin};
  \draw[slflow] (policy) -- node[above,font=\footnotesize]{off} (pin);
  \draw[slflow] (pin) -- node[above,font=\footnotesize]{active} (gate);
  \draw[slflow] (gate.south) -- node[right,font=\footnotesize]{all three agree} (apply.north);
  \draw[slflow] (policy.south) to[bend right=12] node[left,font=\footnotesize]{on} (shadow.north);
  \draw[slflow,dashed] (pin.south) -- node[left,font=\footnotesize]{shadow} (shadow.north east);
\end{tikzpicture}
\caption{The three-tier safety model. A learned weight is applied only when all
three override tiers agree: the operating policy's \texttt{safe\_mode} kill
switch is off, the \texttt{RL\_MODE} deployment pin is \texttt{active}, and the
LB sidecar's gate sees the policy itself return an active action. Any tier
defaulting to its safe state forces \emph{shadow} mode and the data plane falls
back to round-robin. The trained policy alone can never escalate routing.}
\label{fig:dsb:safety}
\end{figure}

\section{Adherence to Codes of Ethics, Regulations, and Laws}
\label{dsb:sec-ethics}

The selected design puts the professional obligations of the ACM/IEEE-CS
Software Engineering Code of Ethics \cite{gotterbarn1999ethics,acm2018ethics}
into the structure of the system itself, rather than bolting them on as policy.
The Code's first principle --- act in the public interest, and hold paramount the
safety of the systems one builds --- is the direct motivation for the
graceful-degradation guarantee. Because every adaptive component degrades to a
deterministic baseline (\cref{dsb:sec-engine}) and three override tiers stand
between any learned decision and live traffic (\cref{dsb:sec-safety}), a defect
or a mis-trained model in the experimental intelligence cannot escalate into an
availability failure for the served application. Safety-by-fallback is the
design's answer to the obligation to minimise the risk the public bears from the
software.

The Code's emphasis on transparency and accountability is met through the audit
log and the operator UI. Every routing and scaling decision is attributable: the
operator can reconstruct, from the durable audit store, why a given backend was
chosen or why the pool changed size at a given moment, and the operator UI shows
what each engine \emph{would} emit even when \texttt{safe\_mode} is suppressing
actuation. The design avoids the opacity of an unexplainable black box: the
deterministic baselines are fully interpretable, and the learned components are
always shadowed by an audited, human-overridable decision path. Our own reporting
follows the same principle --- where a trained engine has not yet met its target
metric, we record that gap plainly rather than paper over it, and the baseline
remains the default until the gap closes.

On data protection, the design's posture is conservative by construction
\cite{gdpr2016}. SmartLoad routes on \emph{operational metadata only} ---
per-backend latency, error rate, request count, and queue depth. It carries no
session affinity and stores no per-client identity (the
stateless-request-path principle), so request bodies and personal data never
enter the decision plane or the telemetry store. The data the system retains
describes the \emph{infrastructure's} behaviour, not the \emph{user's}, which
keeps SmartLoad clear of the personal-data obligations that a content-inspecting
middleware would incur, while still respecting the principle of collecting only
what the function requires.

\section{Response to Global, Economic, Environmental, and Societal Concerns}
\label{dsb:sec-impact}

The four-context concerns raised in Part~A are answered directly by properties of
the selected design rather than by aspiration.

\textbf{Environmental.} The autoscaler is symmetric: \cref{dsb:eq:scale} scales
the pool \emph{down} as readily as up, decommissioning backends within the policy
floor once predicted demand falls. Because the forecaster drives provisioning
proactively, the pool tracks demand rather than being permanently
over-provisioned for a worst case, so idle compute --- and the energy it consumes
--- is released when load recedes. Right-sizing capacity to actual demand is the
design's primary environmental contribution.

\textbf{Economic.} The same scale-down behaviour is a direct cost lever for the
operator, since released instances are instances not billed. The middleware-first
posture adds to this: SmartLoad is a thin layer over commodity NGINX and an
open-source telemetry stack, deployable on the customer's existing infrastructure
with Docker \cite{merkel2014docker}, so adopting it does not mean replatforming
onto a heavyweight mesh or a managed cloud product. Its own overhead stays
bounded because the intelligence runs off the request path and adds no
per-request computation.

\textbf{Global / portability.} The design treats the load balancer as an adapter
behind a stable interface, not a hard dependency: the decision plane speaks to a
generic load-balancer contract, and the NGINX-specific actuation is wrapped up so
that other data planes (for example HAProxy \cite{haproxy} or Envoy
\cite{envoy}) can drop in beside it. In the same spirit, the contract assumes
message loss is possible and the deployment is portable from Docker Compose to an
orchestrator \cite{merkel2014docker} without redesign. This keeps SmartLoad
usable across the varied environments a global user base actually runs.

\textbf{Societal / reliability.} For the people who depend on the served
applications, the most important societal property is that the system stays up.
The layered safety model (\cref{dsb:sec-safety}) and the deterministic floor
(\cref{dsb:sec-engine}) mean that adding adaptive intelligence does not add a new
way for the service to fail: in the worst case SmartLoad behaves exactly like the
classical load balancer it wraps. Reliability is therefore not traded away for
adaptivity --- it is the precondition the adaptivity is built on top of.
